\documentclass[12pt,openright,oneside,a4paper,english,french,spanish]{abntex2}

\usepackage{cmap}
\usepackage{lmodern}
\usepackage[T1]{fontenc}
\usepackage[utf8]{inputenc}
\usepackage{lastpage}
\usepackage{indentfirst}
\usepackage{color}
\usepackage{graphicx}
\usepackage{units}
\usepackage[brazilian,hyperpageref]{backref}
\usepackage[alf]{abntex2cite}
\usepackage{bold-extra}
\usepackage{eso-pic}
\usepackage[dvipsnames,table,xcdraw]{xcolor}
\usepackage{multicol}

% colored hyperlinks
\newcommand{\chref}[2]{%
  \href{#1}{{\usebeamercolor[bg]{Aalborg}#2}}%
}
\setbeamerfont{caption}{size=\tiny}
\setbeamercolor{postit}{fg=white,bg=beamer@headercolor}
\setbeamercolor{alert}{fg=white,bg=red}

\title[]% optional, use only with long paper titles
{Análise da Ferramenta Subversion}

\subtitle{Um sistema de versionamento}  % could also be a conference name

\date{\today}

\author[] % optional, use only with lots of authors
{
  Emilie Morais\\
  Hugo Martins\\
  Ítalo Paiva\\
  Wilton Rodrigues\\
}
% - Give the names in the same order as they appear in the paper.
% - Use the \inst{?} command only if the authors have different
%   affiliation. See the beamer manual for an example

\institute[
 {\includegraphics[width=1cm]{figuras/logo-unb}}\\ %insert a company, department or university logo
  Engenharia de Software\\
  Faculdade\ do Gama\\
  Universidade de Brasília
] % optional - is placed in the bottom of the sidebar on every slide
{% is placed on the bottom of the title page
  Engenharia de Software\\
  Faculdade do Gama\\
  Universidade de Brasília

  %there must be an empty line above this line - otherwise some unwanted space is added between the university and the country (I do not know why;( )
}

% specify the logo in the top right/left of the slide
% \pgfdeclareimage[height=1cm]{mainlogo}{figuras/logo} % placed in the upper left/right corner
% \logo{\pgfuseimage{mainlogo}}

% specify a logo on the titlepage (you can specify additional logos an include them in
% institute command below
\pgfdeclareimage[height=1.5cm]{titlepagelogo}{figuras/logo-unb} % placed on the title page
% \pgfdeclareimage[height=1.5cm]{titlepagelogo2}{figuras/logo-fga} % placed on the title page
\titlegraphic{% is placed on the bottom of the title page
  \pgfuseimage{titlepagelogo}
 % \hspace{5cm}\pgfuseimage{titlepagelogo}
}


% ---
% Início do documento
% ---
\begin{document}
% the titlepage
{\aauwavesbg
\begin{frame}[plain,noframenumbering] % the plain option removes the sidebar and header from the title page
  \titlepage
\end{frame}}
%%%%%%%%%%%%%%%%
\input{fixos/agenda}

% ------------------------------------------------------------
% Inserir slides aqui
% ------------------------------------------------------------
\chapter[Introdução]{Introdução}

  O Subversion, também conhecido apenas como SVN, é um sistema de controle de versão de código aberto. O SVN gerencia arquivos e diretórios, e as modificações feitas neles ao longo do tempo \cite{svn-book}.
  A partir deste controle de versionamento é possível que sejam recuperadas e examinadas as diversas versões de um Item de Configuração, por exemplo.

\section{Histórico}

  O SVN surgiu como um substituto para o Concurrent Versions System (CVS) \cite{cvs-book} no começo do ano 2000. Uma vez que o CVS já havia se firmado como um padrão de sistema de versionamento do mundo open-source não havia nenhum outro sistema de versionamento tão bom quanto o CVS. Por isso a CollabNet, Inc. decidiu então desenvolver um novo sistema de controle de versão do zero \cite{svn-book}.
  O SVN logo atraiu diversos desenvolvedores interessados em contribuir com o projeto, pois muitas pessoas tinham as mesmas dificuldades e experiências frustantes com alguns problemas que eram encontrados no CVS e viram no SVN uma oportunidade de fazer algo a respeito disso.
  Algumas das preocupações da equipe de desenvolvimento do SVN era que qualquer usuário pudesse migrar facilmente do CVS para o SVN. E que o SVN fosse compatível com as características do CVS, embora sem repetir as falhas e bugs dos quais eles queriam tanto se livrar.
  Em 31 de Agosto de 2001 \cite{svn-book}, após 14 meses de desenvolvimento os responsáveis pelo Subversion pararam de usar o CVS para versionar o código fonte e começaram a utilizar o próprio SVN para isso.

\section{Funções gerais}

  A seguir serão descritas as funções que o SVN proporciona:

  \begin{itemize}
  \item Versionamento de diretórios
  Diferente do CVS que apenas rastreia o histórico de arquivos individuais o SVN possui um sistema de arquivos que rastreia as modificações em toda a árvore de diretórios ao longo do tempo. Dessa forma tanto os arquivos quanto os diretórios são versionados.

  \item Histórico de versões efetivo

  Como o CVS é limitado apenas ao versionamento de arquivos, operações como cópia e renomeação—que podem ocorrer com arquivos também, mas que são realmente alterações no conteúdo de algum diretório continente—não são suportadas no CVS. Adicionalmente, no CVS você não pode substituir um arquivo versionado por alguma outra coisa com o mesmo nome sem que o novo item deixe de herdar o histórico do arquivo antigo—que talvez seja até algo com o qual não mantenha nenhuma correlação. Com o Subversion, você pode adicionar, excluir, copiar, e renomear ambos os arquivos ou diretórios. E cada novo arquivo adicionado começa com um histórico próprio e completamente novo.

  \item Commits atômicos

  Um conjunto de modificações ou é inteiramente registrado no repositório, ou não é registrado de forma nenhuma. Isto possibilita aos desenvolvedores criarem e registrarem alterações como blocos lógicos, e também evita problemas que possam ocorrer quando apenas uma parte de um conjunto de alterações seja enviada com sucesso ao repositório.

  \item Versionamento de metadados

  Cada arquivo e diretório tem um conjunto de propriedades—chaves e seus valores—associados consigo. Você pode criar e armazenar quaisquer pares chave/valor que quiser. As propriedades são versionadas ao longo do tempo, tal como os conteúdos de arquivo.

  \item Escolha das camadas de rede

  O Subversion tem uma noção abstrata do acesso ao repositório, tornando-o mais fácil para as pessoas implementarem novos mecanismos de rede. O Subversion pode se associar ao servidor Apache HTTP como um módulo de extensão. Isto dá ao Subversion uma grande vantagem em estabilidade e interoperabilidade, além de acesso instantâneo aos recursos existentes oferecidos por este servidor—autenticação, autorização, compactação online, dentre outros. Um servidor Subversion mais leve e independente também está disponível. Este servidor utiliza um protocolo específico o qual pode ser facilmente ser tunelado sobre SSH.

  \item Manipulação consistente de dados

  O Subversion exprime as diferenças de arquivo usando um algoritmo diferenciado, o qual funciona de maneira idêntica tanto em arquivos texto (compreensível para humanos) quanto em arquivos binários (incompreensível para humanos). Ambos os tipos de arquivos são igualmente armazenados de forma compactada no repositório, e as diferenças são enviadas em ambas as direções pela rede.

  \item Ramificações e rotulagem eficiente

  O custo de se fazer ramificações (branching) e de rotulagem (tagging) não precisa ser proporcional ao tamanho do projeto. O Subversion cria ramos e rótulos simplesmente copiando o projeto, usando um mecanismo semelhante a um hard-link. Assim essas operações levam apenas uma pequena e constante quantidade de tempo.

  \item Hackability

  O Subversion não tem qualquer bagagem histórica; ele é implementado como um conjunto de bibliotecas C compartilhadas com APIs bem definidas. Isto torna o Subversion extremamente manutenível e usável por outras aplicações e linguagens.

  \end{itemize}

\section{Fornecedor}

  Como já mencionado, o SVN foi fundado em 2000 pela CollabNet, Inc., mas agora é desenvolvido como um projeto da Apache Software Foundation \cite{svn-site-home}. Além de receber ajuda de desenvolvedores da comunidade de software livre.

\section{Versões disponíveis}

  Atualmente o site do Subversion recomenda a utilização, apenas, de duas versões, as quais são entituladas de versões suportadas, do SVN \cite{svn-site}. Estas versões suportadas do Subversion no momento são:

  \begin{itemize}
    \item Subversion 1.9.2 (https://subversion.apache.org/download.cgi)

    \item Subversion 1.8.14 (https://subversion.apache.org/download.cgi)

  \end{itemize}

  Embora, se for necessário ou de interesse do usuário utilizar uma versão mais antiga estão disponíveis no link (http://archive.apache.org/dist/subversion/). Neste link é possivel encontrar desde a versão 0.19.1, que é a mais antiga versão disponibilizada, até as versões suportadas pelo Apache Org.

\section{Compatibilidade nos diversos sistemas operacionais}

  O SVN é compatível em todos os Sistemas Operacionais modernos de Unix, Win32, BeOS, OS/2 e MacOS X. O Subversion é escrito em ANSI C e usa APR, a biblioteca Apache Portable Runtime, como uma camada de portabilidade. O cliente Subversion irá funcionar em qualquer lugar que o APR possa ser executado, que é a maioria dos lugares \cite{apache-faq}. O servidor Subversion (ou seja, o lado repositório) é o mesmo, exceto que ele não vai hospedar um repositório Berkeley DB em plataformas Win9x (Win95/Win98/WinME), pois Berkeley DB tem problemas de segmento de memória compartilhada no Win9x. Repositórios FSFS (introduzido na versão 1.1) não tem essa restrição; no entanto, devido a uma limitação no suporte de bloqueio de arquivo do Win9x, eles também não funcionam no Win9x.

  Para reiterar, o cliente Subversion pode ser executado em qualquer plataforma onde APR é executado. O servidor Subversion também pode ser executado em qualquer plataforma onde APR rode, mas não pode hospedar um repositório em Win95/Win98/WinMe \cite{apache-faq}.

\section{Outras informações gerais sobre a ferramenta}

\chapter[Guia de Instalação]{Guia de Instalação}
  
  Esse capítulo apresenta um guia para a instalação do SVN mostrando como instalar, a documentação disponível e as
  principais dificuldades encontradas no processo de instalação.
    

\section{Instalação da ferramenta}
    
    O jeito mais fácil de instalar o Subversion em sistemas Unix-like, como o Ubuntu,
    é utilizando o sistema de pacotes nativo do sistema. Portanto, para a instalação do Subversion,
    utilize o gerenciador de pacotes do Ubuntu \textit{apt-get}, executando o seguinte comando:
    
    \colorbox{Gray}{
      \begin{minipage}{0.6\linewidth}
      \flushleft{ 
      \textbf{sudo apt-get install subversion libapache2-svn}
      }
      \end{minipage}
    }
    
    Com este comando, obtém-se também o pacote de configuração do servidor Apache para o uso com o Subversion.
  
\section{Disponibilidade de documentação de instalação}
  
  A Fundação Apache mantém uma vasta documentação para a instalação do Subversion de diferentes maneiras
  para diferentes sistemas operacionais \footnotemark, porém a documentação não é muito clara e o Subversion necessita
  \footnotetext{Disponível em: \url{http://svn.apache.org/repos/asf/subversion/trunk/INSTALL}.}
  de algumas bibliotecas auxiliares para o correto funcionamento \cite{svn-book}, o que pode complicar demais o processo
  de instalação manual utilizando o código fonte. Portanto, para a instalação do Subversion no Ubuntu, foi utilizado o
  tutorial fornecido em \citeonline{wiki-svn}.
  
  
\section{Principais dificuldades encontradas}
  
  A maior dificuldade encontrada foi na tentativa de instalar o Subversion utilizando o código fonte, pois dependia de 
  várias bibliotecas auxiliares e, mesmo depois de obter essas bibliotecas, não foi possível concluir a instalação da
  ferramenta.
\section{Principais Funcionalidades}

\begin{frame}{Principais Funcionalidades}
\begin{block}{Criar repositórios para manter projetoss}
  \begin{itemize}
    \item<1-> svnadmin create
    \item<1->svn mkdir
    \item<1->svn checkout
  \end{itemize}
\end{block}
\end{frame}
%%%%%%%%%%%%%%%%

\begin{frame}{Principais Funcionalidades}
\begin{block}{Adicionar arquivos para o controle de versão}
  \begin{itemize}
    \item<1-> svnadmin add
  \end{itemize}
\end{block}
\end{frame}
%%%%%%%%%%%%%%%%

\begin{frame}{Principais Funcionalidades}
\begin{block}{Registrar mudanças em um ou mais arquivos}
  \begin{itemize}
    \item<1-> svn commit
    \item<1-> svn revert
    \item<1-> svn update
  \end{itemize}
\end{block}
\end{frame}
%%%%%%%%%%%%%%%%

\begin{frame}{Principais Funcionalidades}
\begin{block}{Rastrear alterações em  um ou mais arquivos}
  \begin{itemize}
    \item<1-> svn status
    \item<1-> svn diff
    \item<1-> svn log
  \end{itemize}
\end{block}
\end{frame}
%%%%%%%%%%%%%%%%

\input{conteudo/verificacao_funcionalidades}
\chapter[Considerações Finais]{Considerações Finais}

\section{Resumo de cobertura das atividades relacionadas às funções da Gerência de Configuração}
 \subsection{Identificação da Configuração}
 \begin{itemize}
 \item \textbf{Requisitos de identificação de dados}
 \item \textbf{Identificar itens de configuação}
 \item \textbf{Estabelecer esquema de identificação}
 \item \textbf{Identificar requisitos de aceitação}
 \item \textbf{Definição das baselines}
 \end{itemize}
 \subsection{Controle de Mudanças}
 \begin{itemize}
 \item \textbf{Estabelecer revisões de controle}
 \item \textbf{EStabelecer procedimento de controle de mudança}
 \item \textbf{Controlar revisões}
 \end{itemize}
 \subsection{Auditoria}
 \begin{itemize}
 \item \textbf{Auditoria física}
 \item \textbf{Auditoria funcional}
 \item \textbf{Revisão formal}
 \end{itemize}
 \subsection{Relato de Estado}
 \begin{itemize}
 \item \textbf{Manter registro da definição dos produtos}
 \item \textbf{Manter registro da verificação da configuração}
 \item \textbf{Manter registro do estado das mudanças}
 \item \textbf{Manter histórico das aprovações das mudanças}
 \end{itemize}
\section{Resumo dos principais achados relacionados à instalação e uso da ferramenta}

	O processo de instalação utilizado para o Ubuntu é bastante simples, embora o processo de instalação manual por meio do 
	código fonte seja bem mais custoso, por depender de outras bibliotecas auxiliares que as vezes não estão inclusas.
	A ferramenta possui uma boa documentação tanto quanto para a instalação, quanto para a descrição das funcionalidades.
	O uso da ferramenta não é tão fácil no primeiro contato, principalmente para quem já está acostumado com outra ferramenta 
	de controle de versão como o Git, pois embora algumas funcionalidades sejam as mesmas, a estrutura de funcionamento difere muito.

	Destacamos abaixo os pontos positivos e negativos em relação à ferramenta.

	\subsection{Pontos positivos}

			Os pontos positivos identificados pelo grupo foram:

			\begin{itemize}
				\item Possui controle de acesso ao repositório;
				\item Gera um log com data e nome do autor que fez aquela alteração;
				\item Trabalha com a criação de \textit{branches};
			\end{itemize}


	\subsection{Pontos negativos}

			Os pontos negativos identificados pelo grupo foram:

			\begin{itemize}
				\item Necessidade do uso de um servidor próprio;
				\item Muitos passos para configurar;
				\item Não há possibilidade de realizar \textit{commits} locais;
			\end{itemize}

\section{O grupo usaria a ferramenta em seus projetos de desenvolvimento de software? Por quê?}


% ----------------- Referências --------------------------------

% --- O comando \allowframebreaks ---
% Se o conteúdo não se encaixa em um quadro, a opção allowframebreaks instrui
% beamer para quebrá-lo automaticamente entre dois ou mais quadros,
% mantendo o frametitle do primeiro quadro (dado como argumento) e acrescentando
% um número romano ou algo parecido na continuação.

\begin{frame}{Referências}
  \nocite{*}
  \bibliographystyle{plain}
  \bibliography{editaveis/bibliografia}
\end{frame}

% ----------------- FIM DO DOCUMENTO -----------------------------------------

{\aauwavesbg%
\begin{frame}[plain,noframenumbering]%
  \finalpage{Obrigado!}
\end{frame}}
%%%%%%%%%%%%%%%%

\end{document}
